% Auxiliary degree estimate. It does not supply a hereditary object atlas.
\begin{proposition}[Proved: a degree-linear exceptional core]
\label{D:exceptional-degree-linear-core}
Use the fixed surface resolution, strictly anti-nef cycle
$A=\sum_{i=1}^{s}a_iE_i$, and strict original boundary $D$ of
Proposition~\ref{D:fixed-boundary-uniform-different}.  Put
\[
 a=\sum_i a_i,\qquad
 b=(K_Y+D)\cdot A+a,\qquad c=3b+a.
\]
Let a connected finite cover of the indicated open have degree
$n$, and assume it is tame at the original boundary divisors.
The cover need not be Galois.  Normalize $Y$ in its function
field, resolve, and contract all smooth rational $(-1)$-curves
which map to points of $Y$, as in the proof of that proposition.
Write $f:Z\to Y$ for the resulting morphism and
$C=(f^*A)_{\mathrm{red}}$.

Define the \emph{core} $\mathcal K$ to consist of every component
of $C$ dominating an $E_i$, together with every point-mapped
component $Q$ for which $K_Z\cdot Q>0$.  Then
\[
 \#\mathcal K\le n(b+a),\qquad
 \sum_{F\in\mathcal K}(-F^2)\le nc,\qquad
 \sum_{F\in\mathcal K}p_a(F)\le nb.                 \tag{D.GC.1}
\]
Every component outside the core is a smooth rational
$(-2)$-curve.  If $L$ is the sum, with intersection
multiplicity, of all incidences from a core component to a
noncore component, then
\[
                 L<2nc\quad\text{if }\mathcal K\ne\varnothing,
                 \qquad L=0\quad\text{otherwise}.       \tag{D.GC.2}
\]
The sum of intersection multiplicities over unordered pairs of
distinct core components is less than $nc/2$ when the core is nonempty.
There are at most $2nc+n$ connected noncore blocks.
The numbers of vertices within the negative-definite rational
$(-2)$-blocks remain unrestricted by these estimates.
\end{proposition}

\begin{proof}
Write the components dominating $E_i$ as $F_{i\lambda}$ and
the other components of $C$ as $Q_j$.  Their coefficients in
$f^*A$ are $m_{i\lambda}=a_i e_{i\lambda}$ and positive
integers $m_j$, respectively.  The positivity of $m_j$ follows
because $A$ is an effective Cartier divisor whose support
contains the image of $Q_j$.  The degree formula at each base
divisor gives
\[
 \sum_\lambda e_{i\lambda}f_{i\lambda}=n,
 \qquad
 \sum_{i,\lambda}m_{i\lambda}\le na.                    \tag{D.GC.3}
\]
Here the residue degrees are full residue degrees; no
separability of residual extensions is assumed.

The determinant calculation and horizontal tame correction in
Propositions~\ref{D:surface-uniform-different} and
\ref{D:fixed-boundary-uniform-different} give, before any use
of Galois conjugacy,
\[
                 K_Z\cdot f^*A\le n(K_Y+D)\cdot A.
\]
Adjunction and the chosen contractions give
$K_Z\cdot Q_j\ge0$ and $K_Z\cdot F_{i\lambda}\ge-1$.
Consequently the following sum has nonnegative summands:
\[
 \sum_{i,\lambda}m_{i\lambda}
                 (K_Z\cdot F_{i\lambda}+1)
       +\sum_jm_jK_Z\cdot Q_j\le nb.                  \tag{D.GC.4}
\]
In particular $b\ge0$.  There are at most $na$ dominating
components.  Each point-mapped core component contributes at
least one to the second sum in (D.GC.4), proving the first
bound in (D.GC.1).

All these curves are exceptional over points of the normal
cover of $S$, so their intersection matrix is negative
definite and each square is at most $-1$.  For a dominating
component $F$, adjunction gives
\[
 -F^2\le K_Z\cdot F+2,
 \qquad p_a(F)\le (K_Z\cdot F+1)/2.
\]
For a point-mapped core component $Q$, its positive integral
canonical degree gives
\[
 -Q^2\le K_Z\cdot Q+2\le3K_Z\cdot Q,
 \qquad p_a(Q)\le(K_Z\cdot Q+1)/2\le K_Z\cdot Q.
\]
Apply (D.GC.3)--(D.GC.4), using all $m\ge1$, to obtain the
remaining bounds in (D.GC.1).  A noncore point-mapped curve
satisfies
$0=2p_a(Q)-2-Q^2$.  Its square is negative; hence its arithmetic
genus is zero and its square is $-2$.  An integral proper
curve of arithmetic genus zero over an algebraically closed
field is a smooth rational curve.

We give the attachment estimate directly, without needing a
classification of the possible $(-2)$-graphs.  In the negative
intersection matrix, order the core first and the noncore
blocks second, and write
\[
                   M=\begin{pmatrix}M_{\mathcal K}&-T\\
                                      -T^{\mathsf t}&N
                      \end{pmatrix}.
\]
All entries of $T$ are nonnegative integers.  Each diagonal
entry of $N$ is $2$, its off-diagonal entries are nonpositive,
and $N$ is positive definite.  Writing $N=2I-P$ with $P$
symmetric and entrywise nonnegative, positive definiteness
gives $\rho(P)<2$.  Thus
\[
             N^{-1}=\frac12\sum_{j\ge0}(P/2)^j
\]
is entrywise nonnegative and has diagonal entries at least
$1/2$.  The Schur complement
$M_{\mathcal K}-TN^{-1}T^{\mathsf t}$ is positive definite.
For the row $t_F$ corresponding to a core component $F$,
\[
     -F^2>t_FN^{-1}t_F^{\mathsf t}
             \ge\tfrac12\sum_Q(F\cdot Q)^2
             \ge\tfrac12\sum_Q F\cdot Q.
\]
Summing proves (D.GC.2).  Applying positive definiteness of
$M_{\mathcal K}$ to its all-one vector gives the stated bound
on core--core incidences.

Every noncore block in a connected component meeting the
core has at least one attachment, so their number is at most
$L$.  A connected component not meeting the core is one
noncore block.  The whole exceptional fibre has at most $n$
connected components: the finite normal cover has at most
$n$ points over $s$, and proper birational maps to its normal
local germs have connected fibres.  This proves the last
bound.  The same argument applies after strict henselization,
where the connected normalization has only one closed point.
\end{proof}

\begin{remark}[Remaining presentation problem]
These are numerical and intersection-graph bounds.  We do
not assert that contracting the noncore blocks yields an
object atlas with $O(nr^2)$ parameters, or that the constant
in the quadratic representation bound grows linearly with
cover degree.  A uniform treatment of the full module and
Frobenius data along arbitrarily long $(-2)$-blocks remains
necessary for that argument.
\end{remark}
